% ======================================================================
% Section 8 – Livrable spécifique : entente de gouvernance des données simplifiée
% Grille : « Livrable spécifique » – 10 pts (section prioritaire).
% Longueur : 2 à 3 pages (vérifier avec « make entente » → build/entente.pdf).
%
% Ce fichier est inclus DANS LE RAPPORT et DANS rapport/entente-autonome.tex :
% ne rien y écrire qui n'ait de sens que dans le rapport, sauf dans le
% bloc \ifententeautonome\else … \fi ci-dessous.
%
% Exigences : entente entre le consortium et UNE communauté participante,
% claire, applicable et cohérente avec les principes PCAP. Articles
% numérotés avec \articleentente{Titre}.
% ======================================================================
\setcounter{articleentente}{0}

\ifententeautonome\else
\section{Livrable spécifique : entente de gouvernance des données simplifiée}
\label{sec:entente}
% TODO : un court paragraphe de présentation (rapport seulement) : rôle de l'entente, lien avec les sections 5 et 6.
\arediger{présentation de l'entente}
\fi

\articleentente{Parties}
\label{art:parties}
% TODO : identifier le consortium \fedesante{} (représentant) et la communauté participante (instance de gouvernance, centre de santé).
\arediger{parties}

\articleentente{Objet}
\label{art:objet}
% TODO : objet de l'entente (participation à l'entraînement fédéré, portée, exclusions).
\arediger{objet}

\articleentente{Définitions}
\label{art:definitions}
% TODO : définir au minimum : données locales, mise à jour du modèle, modèle global, \acs{DP}, désapprentissage, publication, bénéfices.
\arediger{définitions}

\articleentente{Droits PCAP de la communauté}
\label{art:droits-pcap}
% TODO : traduire chacun des quatre principes (propriété, contrôle, accès, possession) en droits concrets et vérifiables.
\arediger{droits PCAP}

\articleentente{Conditions de retrait et désapprentissage}
\label{art:retrait}
% TODO : préavis, effets du retrait, engagement de désapprentissage (ou de réentraînement sans la contribution), délais, preuve fournie.
\arediger{retrait et désapprentissage}

\articleentente{Approbation des publications}
\label{art:publications}
% TODO : processus d'approbation préalable des publications et présentations par la communauté ; délais ; droit de refus.
\arediger{approbation des publications}

\articleentente{Partage des bénéfices}
\label{art:benefices}
% TODO : bénéfices (accès au modèle, retombées financières, renforcement des capacités, reconnaissance).
\arediger{partage des bénéfices}

\articleentente{Durée}
\label{art:duree}
% TODO : durée, renouvellement, sort du modèle et des documents à l'échéance.
\arediger{durée}

\articleentente{Résolution des différends}
\label{art:differends}
% TODO : étapes (discussion, médiation, instance désignée) ; respect des modes de gouvernance de la communauté.
\arediger{résolution des différends}

\articleentente{Signatures}
\label{art:signatures}
% TODO : vérifier les titres des signataires.
\vspace{0.5em}
\noindent
\begin{tabularx}{\linewidth}{@{}XX@{}}
Pour le consortium \fedesante{} & Pour la communauté participante \\[2.5em]
\rule{0.9\linewidth}{0.4pt} & \rule{0.9\linewidth}{0.4pt} \\
Nom, titre : [\ldots] & Nom, titre : [\ldots] \\
Date : [\ldots] & Date : [\ldots] \\
\end{tabularx}
